\section{Negative time: smoothing and the completed fluctuation hierarchy}
\label{sec:negative}

\subsection{The remote negative endpoint}
Let $\phi_0=\Phi(0)>0$, $C=\phi_0\sqrt\pi/2$, and
\begin{equation}
 N_T(z)=C T^{-1/2}e^{-z^2/(4T)}.
 \label{eq:negative-gaussian}
\end{equation}
An explicit convergent source-series majorant supplies
$B_{2,\Phi}\ge\sup_{u\ge0}|\Phi''(u)|$. Evenness gives
$|\Phi(u)-\phi_0|\le B_{2,\Phi}u^2/2$. Substitution in
\eqref{eq:source}, followed by $v=\sqrt T u$, proves
\begin{equation}
 |\partial_z^j(H(-T,z)-N_T(z))|
 \le\frac{B_{2,\Phi}}2 T^{-(j+3)/2}J_{j+2}(A),
 \quad |y|\le A\sqrt T,
 \label{eq:negative-absolute}
\end{equation}
for each fixed derivative order $j$, uniformly in real $x$. Here
\begin{align*}
 J_m(A)&=\int_0^\infty v^m e^{-v^2+Av}\dd v,\\
 J_0(A)&=\frac{\sqrt\pi}2e^{A^2/4}\operatorname{erfc}(-A/2),
 &J_1(A)&=\frac{1+AJ_0(A)}2,\\
 J_{m+1}(A)&=\frac{mJ_{m-1}(A)+AJ_m(A)}2\quad(m\ge1).
\end{align*}
This elementary proof also explains the normalization of the Gaussian.
It is an absolute estimate. At very large heights the Gaussian can be
far smaller than its error; no uniform relative high-height asymptotic
is asserted.

For a fixed constant regulator $\epsilon>0$, the unaveraged density
defined in Section~\ref{sec:physical} has the scaled endpoint profile
\begin{equation}
 T^2\rho_{\epsilon}(-T,\sqrt T(\xi+\ii\upsilon))
 =\frac{\phi_0^2}{16\epsilon^2}(\xi^2+\upsilon^2)
       e^{-(\xi^2-\upsilon^2)/2}+O(T^{-1})
 \label{eq:negative-density-profile}
\end{equation}
on compact $(\xi,\upsilon)$ sets. The physical adaptive regulator used
later is different; its endpoint transfer has its own error rates.
The parent endpoint proof gives both accounts \cite{parent}.

\subsection{Frozen high-height charts}
Fix $0<\tau\le T\le B$, observation widths $r,b$, and a large real
center $X$. Write
\[
 L=\log(X/(4\pi)),\quad N=\lfloor\sqrt{X/(4\pi)}\rfloor,
 \quad V_B=8\sqrt{\max(1,B)}L,\quad d=z-X.
\]
All of these choices are frozen on the chart. Put
\[
 w_\epsilon(z)=\frac{1+\epsilon\ii z}2,\quad
 \ell_\epsilon(z)=\log M(w_\epsilon(z)),\quad
 G(z)=\sqrt{M(w_+(z))M(w_-(z))},\qquad \epsilon\in\{+1,-1\},
\]
with compatible analytic branches and $G(X)>0$. On fixed buffers,
$|G(z)|$ has the scale $X^{7/4}e^{-\pi X/8}$ up to bounded ratios.
The gamma norm is
\begin{equation}
 \|E\|_\gamma^{(3,1)}=
 \sup_{\substack{\tau\le T\le B\\|\Re z-X|\le r,\ |\Im z|\le b}}
 \frac{\max_{0\le j\le3}|\partial_z^j E|+|\partial_T E|}{|G(z)|}.
 \label{eq:gamma-norm}
\end{equation}

For each arm and arithmetic mode set
\begin{align*}
 b_{\epsilon n}&=M(w_\epsilon(X))n^{-w_\epsilon(X)},\\
 \lambda_{\epsilon n}&=\ell_\epsilon'(X)-\epsilon\ii\log n/2,
 &c_{\epsilon r}&=\frac{\ell_\epsilon^{(r+1)}(X)}{(r+1)!}.
\end{align*}
The fluctuation factor is
\[
 \exp\{\ell_\epsilon(X+q)-\ell_\epsilon(X)-\ell_\epsilon'(X)q\}.
\]
Its finite weighted polynomials are \eqref{eq:finite-partition} with
$d=1$ and $c_r=c_{\epsilon r}$. A simultaneous Cauchy estimate on a
large disk gives $|c_{\epsilon r}|\le A D^{r-1}X^{-r}$ for every needed
index. This common bound, rather than separate unspecified constants
at each order, makes controlled order growth possible.

Gaussian evolution of a polynomial is explicit. With
\begin{equation}
 \mathcal H_m(u,T)=m!\sum_{\nu=0}^{\lfloor m/2\rfloor}
                 \frac{T^\nu u^{m-2\nu}}{\nu!(m-2\nu)!},
 \quad \mathcal H_m=u\mathcal H_{m-1}+2T(m-1)\mathcal H_{m-2},
 \label{eq:heat-moments}
\end{equation}
the heat transform of $q^m e^{\lambda q}$ is
$e^{\lambda d+T\lambda^2}\mathcal H_m(d+2T\lambda,T)$.
Equivalently, define $Q_T=d+2T\partial_d$ and
\begin{align}
 U_{\epsilon0}&=\sum_{n\le N}b_{\epsilon n}
                          e^{\lambda_{\epsilon n}d+T\lambda_{\epsilon n}^2},\\
 jU_{\epsilon j}&=\sum_{r=1}^j r c_{\epsilon r}
                            Q_T^{r+1}U_{\epsilon,j-r}.
 \label{eq:negative-recurrence}
\end{align}
Thus every retained grade is a finite Gaussian evolution of fixed real
initial data and solves the heat equation exactly in $T$.

\subsection{The arithmetic completion is part of the model}
A bare frozen sum misses cutoff crossings and the Riemann--Siegel
remainder. On $[X-V_B,X+V_B]$ let
\[
 a(u)=\sqrt{u/(4\pi)},\quad m(u)=\lfloor a(u)\rfloor,
 \quad\chi(u)=1-2(a(u)-m(u)),
\]
and let $S_m(u)=\sum_{\epsilon,n\le m}M(w_\epsilon(u))n^{-w_\epsilon(u)}$.
Use the fixed real Riemann--Siegel correction
\begin{align*}
 \mathcal U(u)&=\exp\{-\ii[u\log(u/(4\pi))/4-u/4-\pi/8]\},\\
 R_{m,K}(u)&=2\Re\left\{(-1)^{m-1}M(w_+(u))\mathcal U(u)a(u)^{-1/2}
       \sum_{k=0}^K C_k^{\rm RS}(\chi(u),1/2)a(u)^{-k}\right\}.
\end{align*}
The coefficient functions $C_k^{\rm RS}$ are the fixed canonical
functions in the supporting recurrence manuscript, with its finite
providers and removable-value conventions. Define
\begin{align}
 I_{X,K,B}(T,z)&=\int_{X-V_B}^{X+V_B}
 \frac{e^{-(z-u)^2/(4T)}}{\sqrt{4\pi T}}
       [S_{m(u)}(u)-S_N(u)+R_{m(u),K}(u)]\dd u,\\
 F^{[p,K]}(T,z)&=\sum_{\epsilon=\pm1}\sum_{j=0}^p U_{\epsilon j}(T,z)
                     +I_{X,K,B}(T,z).
 \label{eq:negative-completed}
\end{align}
The integrand is evaluated as fixed real data on its fixed partition.
No moving arithmetic floor is differentiated. Consequently
$(\partial_T-\partial_{zz})F^{[p,K]}=0$ exactly. Finite-precision outputs
have a separate evaluation error; their rounded values are not themselves
claimed to satisfy an exact differential equation.

\begin{theorem}[Completed fixed order; recurrence manuscript]
\label{thm:negative-fixed}
For fixed $p\ge1$, $0<\tau\le B$ and fixed observation buffers, there
are finite constants and a sufficient height onset such that
\begin{equation}
 \|H(-T,\cdot)-F^{[p,2p]}(T,\cdot)\|_\gamma^{(3,1)}
          \le C_{p,r,b}(\tau,B)X^{-p-3/4}.
 \label{eq:negative-fixed-error}
\end{equation}
\end{theorem}
\begin{proof}[Proof mechanism]
The local weighted-exponential remainder is paid with its Gaussian
moments. Summing the absolute arithmetic population costs $X^{1/4}$.
The canonical Riemann--Siegel initial remainder is bounded by
\[
 \frac{2^{7/4}}7G(u)a(u)^{-1/2}
       \Gamma((K+1)/2)\left(\frac{1.1}{a(u)}\right)^{K+1},
\]
so $K=2p$ gives the same power $X^{-p-3/4}$ after heat propagation.
The cutoff correction pays every jump; the omitted real heat tails
are kept separately. Spatial and time kernel derivatives are bounded
on larger buffers before restriction. This proves the asserted norm
with the explicit finite constant construction in \cite{recurrence}.
\end{proof}
The detailed proof is the recurrence companion's
\src{recurrence/manuscript/sections/04_comparison.tex}{fixed-order theorem}
(\texttt{thm:fixed}). The constants contain inverse powers of $\tau$ and factors such as
$e^{(b+1)^2/(4\tau)+2B}$. The theorem does not allow an unproved
substitution $\tau=\tau(X)$ or $B=B(X)$.

\subsection{Growing order, onset and finite evaluation}
Passing from the gamma envelope to the outer envelope of
Section~\ref{sec:crossover} costs at most a fixed multiple of
$e^{L^2/8}$ on $1\le T\le2$. The simultaneous coefficient estimates
control the remaining order dependence by $e^{O(L\log L)}$.
Keeping the heat tails leads to the following stronger conclusion
\cite{recurrence,rec-onset,rec-bounds}.

\begin{theorem}[Controlled growing order]
\label{thm:negative-growing}
For fixed $a>1/8$, $p=\lfloor aL\rfloor$, and
$0<\eta<\min(a-1/8,127/8)$, the completed error on $1\le T\le2$
is at most $e^{-\eta L^2}$ in the outer version of
\eqref{eq:gamma-norm}, beyond a sufficient onset. In particular
$p=\lfloor L/2\rfloor$ permits $\eta=1/4$.
\end{theorem}

The recurrence companion gives the full growing-order proof
(\src{recurrence/manuscript/sections/05_growing_order.tex}{\texttt{thm:growing}}).
For this specified $a=1/2$ construction, write
$\nu=\lceil r+b+2\rceil+1$. The sharpened sufficient analytic onset is
\begin{equation}
 L\ge64\nu;
 \qquad X\ge2^{128\nu+8}\quad\hbox{is a sufficient dyadic choice}.
 \label{eq:rec-onset}
\end{equation}
The older $L\ge2^{22}\nu^4$ onset remains a valid fallback, not the
preferred requirement. Strict domain and denominator guards remain
part of the certified descriptors.

The finite evaluator returns the field, spatial derivatives through
three and the time derivative with a combined analytic and evaluation
allowance $\Eo e^{-L^2/4}$. For rational query data and dyadic centers,
the delivered precision has scale
\[
 P=\frac{3L^2}{8\log2}+O(L\log L).
\]
With $Q$ the sum of input length, delivered precision, separation bits,
$(\lceil L\rceil)^2$ and a fixed offset, the proved coarse bounds are
\begin{equation}
 \operatorname{work}\le C(N+1)Q^{40},\qquad
 \operatorname{memory}\le C Q^{20}.
 \label{eq:negative-complexity}
\end{equation}
The arithmetic population $N\asymp\sqrt X$ has not disappeared.
The uniform majorants $B_F$ and every $D_k$ are multiplied by
$\operatorname{csch}(\pi)<1/10$. Separately, the enlarged-contour
majorant $B_c/16$ permits eleven fewer Gauss nodes per edge,
forty-four per instantiated four-edge rule, at the same target.
These are analytical contour-count improvements; they do not reduce
the outer completion-integral nodes or establish measured runtime gains.
The exact input guards, half-target onset and delivered evaluation contract
are in the recurrence companion's
\src{recurrence/manuscript/sections/11_calibration_and_onset.tex}{\texttt{thm:linear-onset}}
and \src{recurrence/manuscript/sections/12_effective_evaluation.tex}{\texttt{thm:effective}, \texttt{thm:bit-cost}}. Component replays support the finite provider recipes;
a complete software implementation is a separate task.

\subsection{The integrated moving lower endpoint}
Recurrence Note 20's proof is integrated as
\src{recurrence/manuscript/sections/14_moving_endpoint.tex}{\texttt{thm:moving-endpoint-delivery}},
with its distinct \texttt{def:moving-channel-norm} \cite{rec-moving}.
It delivers field, spatial orders zero through three and direct time
error together in the outer $(3,1)$ norm on $L^{-1}\le T\le2$:
\[
 \|H(-T,\cdot)-\widehat F\|_{\mathrm o,[1/L,2]}^{(3,1)}
 \le e^{-L^2/4},\qquad L\ge1024\nu^2.
\]
The construction keeps $p=\lfloor L/2\rfloor$, $K=2p$,
$V=8\sqrt2L$, all arithmetic modes, exact adjacent-cutoff transfers
and the bounded completion integral. The new proof pays
$e^{Y^2/(4T)}$, the Gaussian normalization, inverse-time powers,
complete population and degree-dependent evolved moments jointly.
Substitution into a fixed-$\tau$ theorem would not supply these payments.

Ordinary rational queries carry a strict witness $T L^->1$ for a
certified rational lower bound $L^-\le L$. The literal $T=1/L$ tag
has a separately paid time-input error channel, including $F_{TT}$.
There is no arbitrary-real equality oracle. Delivered precision remains
$3L^2/(8\log2)+O(L\log L)$, but the unchanged outer ellipse now needs
$O_{r,b}(L^3)$ sufficient completion-quadrature nodes. Internal provider
work and query separation are additional costs. The preserved $[1,2]$
theorem retains its smaller $64\nu$ onset and its own outer-node count.
Neither theorem includes $T=0$ or a complete certified software evaluator.
